\clearpage
\fancyhead[L]{\small BODE. ASTRONOMISCHES JAHRBUCH.}
\hypertarget{native-553-lower}{}\label{native-553-lower}
\begin{center}\small Göttingische gelehrte Anzeigen. 22. Stück. Seiten 209..220. 6. Februar 1813.\\[3pt]\rule{.72\linewidth}{.3pt}\end{center}

\emph{Astronomisches Jahrbuch für das Jahr} 1814, mit einer Sammlung der neuesten Abhandlungen, Beobachtungen und Nachrichten aus den astronomischen Wissenschaften, berechnet und herausgegeben von J. E. \textsc{Bode}, königlichem Astronomen und Miglied der Akademie der Wissenschaften. 1811. Beim Verfasser und in Kommission bei J. E. \textsc{Hitzig}. Berlin. 275 Seiten in Oktav und eine Kupfertafel.

Astronomisches Jahrbuch für das Jahr 1815 usw. Ebenda 1812. Ebenfalls 275 Seiten in Oktav und eine Kupfertafel.

Wir wiederholen die Bemerkungen zur Einrichtung des astronomischen Kalenders nicht, die wir bereits bei der Besprechung früherer Bände dieses Jahrbuchs gemacht haben, und beschränken uns hier auf die Anzeige der beigefügten Abhandlungen. Der Jahrgang 1814 enthält die folgenden 34 Artikel. 1) Die mittlere astronomische Strahlenbrechung nach \textsc{Laplace}. 2) Die astronomischen Beobachtungen, die \textsc{Bugge} in den Jahren 1809 und 1810 an der königlichen Sternwarte in Kopenhagen anstellte, betreffen Ceres, Vesta, Uranus, Saturn, Jupiter, die Verfinsterungen seiner Trabanten und Sternbedeckungen. 3) Über eine Methode zur Zeitbestimmung durch Messung des Abstands der Sonne von einem festen und bekannten Punkt am Horizont, von Herrn J. F. \textsc{von Beek Calkoen}. Dies \SourceLeaf{554}Methode kann bei gebotener Vorsicht brauchbare Ergebnisse liefern. Sie darf aber nicht auf Gegenstände am Horizont beschränkt werden, und Refraktion und Parallaxe dürfen nicht vernachlässigt werden. Beides lässt sich angemessen berücksichtigen, ohne das Verfahren sehr umständlich zu machen. 4) Die astronomischen Beobachtungen, die Herr Doktor \textsc{Triesnecker} und Herr Professor \textsc{Bürg} an der kaiserlichen Sternwarte in Wien anstellten, umfassen Trabantenverfinsterungen, Sternbedeckungen sowie Planeten- und Sonnenbeobachtungen. 5) Die Ideen zur Störungsrechnung nach \textsc{Kepler} samt Anmerkungen von Herrn J. W. \textsc{Pfaff} verdienen zum Teil Beachtung, obwohl uns manches unausgereift oder übereilt erscheint. 6) Astronomische Beobachtungen, 1810 an der königlichen Sternwarte in Prag von Herrn Professor \textsc{David} und Herrn Adjunkt \textsc{Bittner} angestellt. Der Art, wie hier und in den vorangegangenen Bänden dieses Jahrbuchs beobachtete Sternhöhen für eine vermeintliche Korrektur der Refraktionstafeln verwendet werden, können wir nicht zustimmen. Wer über diesen heiklen Punkt arbeiten will, muss zunächst die Polhöhe seines Beobachtungsorts zuverlässiger bestimmen, die Beobachtungen insgesamt weit stärker vermehren, als Herr \textsc{David} es getan hat, und darf sich nicht damit begnügen, die Sterndeklinationen von anderen Astronomen zu übernehmen. 7) Aus einem Schreiben von Herrn Professor \textsc{Littrow} aus Kasan. Es ist erfreulich, daraus zu ersehen, wie die Astronomie an der Grenze Asiens aufblüht. 8) Genauere Bestimmung der Periode der Helligkeitsänderung des Sterns $\eta$ Antinoi, von Herrn Professor \textsc{Wurm}. Aus eigenen Beobachtungen, die einen Zeitraum von 24 Jahren umfassen, findet Herr \textsc{Wurm} die Periode $7{,}1761$ Tage. Man hätte sich hier eine weniger willkürliche Art der Zusammenfassung der Beobachtungen wünschen können, etwa nach der Methode der kleinsten Quadrate, und anschließend einen Vergleich mit den einzelnen Beobachtungen. 9) Beobachtungen des Kometen von 1807 und der totalen Sonnenfinsternis vom 16. Juni 1806 in Salem in den Vereinigten Nordamerikanischen Staaten, von Herrn \textsc{Bowditch}. Aus dem letzten Band der \emph{Memoirs of the American Academy of Arts and Sciences}. Besonders interessant ist die Schilderung der Begleitumstände der totalen Sonnenfinsternis. 10) Nachricht über die Mannheimer Sternwarte von Herrn Staatsrat \textsc{Klüber}, Kurator der Sternwarte. Was hier gesagt wird, ist durch die seither erschienene gesonderte Schrift von Herrn \textsc{Klüber} bereits ausführlicher bekannt. 11) Astronomische Beobachtungen, 1810 an der königlichen Sternwarte in Berlin angestellt. Wir heben daraus die Bestimmung der Polhöhe der Berliner Sternwarte mit dem zweifüßigen Troughton-Kreis hervor, für die Herr \textsc{Bode} im Mittel $52^\circ\,31'\,15''$ findet. 12) Über den Kometen von 1795, von Herrn Dr. \textsc{Olbers}. Die Beobachtungen dieses Kometen sind sehr dürftig, und die von verschiedenen Astronomen ermittelten Elemente weichen beträchtlich voneinander ab. Herr Dr. \textsc{Olbers} stützte eine neue Bestimmung der Elemente teils auf eine zuvor noch nicht verwendete Beobachtung von \textsc{Herschel}, teils auf eine neu reduzierte Beobachtung von Herrn \textsc{Bode}. 13) Ergebnisse einer Untersuchung über die Lage der Ebene des Saturnrings, die Theorie des vierten Satelliten, die Massen des Planeten und des Rings sowie beobachtete Sternbedeckungen, von Herrn Professor \textsc{Bessel}. Die ausführliche Untersuchung zum vierten Saturntrabanten ist bereits im dritten Heft des Königsberger Archivs erschienen. Die wichtigen Ergebnisse verdienen die Aufmerksamkeit aller Astronomen. 14) Berechnung der Bahn des Kometen von 1810, vom selben Verfasser. 15) Astronomische Beobachtungen, die Herr \textsc{Derfflinger} 1809 und 1810 in Kremsmünster anstellte. Beim Eintritt des Aldebaran am 18. September 1810 schien der Stern für Herrn \textsc{Derfflinger} ein paar Sekunden lang mit verändertem Licht am hellen Mondrand zu verweilen. Dasselbe bemerkten auch \textsc{Bugge}, \textsc{Bode} und \textsc{David}, und einige Astronomen haben es für etwas Außerordentliches gehalten. An unserer Sternwarte, wo dieselbe Bedeckung ebenfalls beobachtet wurde, ist nichts dergleichen bemerkt worden. Wir lassen daher offen, ob es sich um etwas anderes als die Irradiation des Mondrands bei vielleicht schon etwas ermüdetem Auge gehandelt hat. 16) Beobachtung und Berechnung der Bedeckung des Aldebaran durch den Mond am 18. September 1810, \SourceLeaf{555}in Dorpat von Herrn Professor \textsc{Knorre} vorgenommen. Von dem eben erwähnten Phänomen ist hier nichts gesagt. 17) Über die Höhenmessung mit dem Barometer, von Herrn Professor \textsc{Benzenberg}. Dieser an sich lobenswerte Versuch, den Gegenstand allgemeinverständlich darzustellen, ist hier wohl nicht ganz am Platz. 18) Beobachtungen von Juno und Vesta im Jahr 1811 an der kaiserlichen Sternwarte in Wilna, von Herrn Professor \textsc{Sniadecki}. 19) Aus einem Schreiben von Herrn Dr. \textsc{Pansner} in Petersburg. Dieser Artikel berichtet über eine trigonometrische Vermessung der Küste des Finnischen Meerbusens, deren Abschluss sehr zu wünschen ist. 20) Über die Genauigkeit des Vertikalkreises von \textsc{Baumann}, von Herrn Dr. \textsc{Pottgiesser} in Eberfeld. Der Aufsatz zeugt von großer praktischer Einsicht. Der Verfasser hat die Fehler anhand der Beschaffenheit des Instruments selbst abgeschätzt und dessen Genauigkeit außerdem anhand einiger tatsächlicher Beobachtungen geprüft, die sehr zu seinen Gunsten sprechen. Wir hätten uns nur gewünscht, dass die Beobachtungen zahlreicher und vielfältiger wären, da neuere Erfahrungen einige Bedenken gegen Kreise mit stehenden Säulen geweckt haben. 21) Bestimmungen der Länge und Breite einiger Orte im österreichischen Gebiet, dazu beobachtete Sternbedeckungen, von Frau Reichsfreiin von \textsc{Matt}. 22) Beobachtungen zur jährlichen Parallaxe von $\alpha$ in der Leier, von Herrn \textsc{Calandrelli} in Rom. Hier wird auf eine jährliche Parallaxe von $5''8$ geschlossen, gegen deren Zuverlässigkeit sich allerdings noch einiges einwenden ließe. Der Sektor, mit dem die Beobachtungen angestellt wurden, wurde nicht umgewendet. Man setzte voraus, dass sein Kollimationsfehler das ganze Jahr über unverändert geblieben sei. 23) Entwurf einer Sonnenuhr, die zwölf Uhr mittags nach mittlerer Zeit anzeigt. Die in Form einer 8 geschlungene Kurve, die die Spitze des Zeigerschattens auf einer vertikalen Sonnenuhr in den verschiedenen Monaten des Jahres am mittleren Mittag beschreibt, wird durch 36 Punkte gezeichnet. Die Verwendung der mittleren Zeit für öffentliche Uhren hat allerdings viel für sich. Da aber diejenigen, die auf dem Land oder in kleineren Orten für das Stellen der Uhren zuständig sind, darüber selten hinreichend unterrichtet sind, geht dabei der große Vorteil der \emph{Übereinstimmung} der öffentlichen Uhren leicht verloren. An dieser Übereinstimmung ist besonders auf Poststraßen viel gelegen, und bei Verwendung der wahren Zeit lässt sie sich leichter erhalten. 24) Astronomische Ortsbestimmungen von Herrn Oberprediger \textsc{Fritsch} in Quedlinburg. Herr \textsc{Fritsch} bestimmte auf einer Reise nach Schlesien im Sommer 1810 die Polhöhe von Ballenstedt, Werningerode, Oschatz, Bischofswerda und mehreren Punkten in Schlesien. 25) Über den Kometen des Jahres 1811 und seine Wiederkehr im August, von Herrn Dr. \textsc{Olbers}. 26) Astronomische Beobachtungen und Bemerkungen von Herrn Professor \textsc{Gauss}. Sie enthalten Beobachtungen von Pallas und Vesta und Berichtigungen einiger Bemerkungen von Herrn Dr. \textsc{Treisnecker} im Jahrbuch für 1813. Die drei folgenden Artikel geben die Ephemeriden für Pallas, Juno und Vesta von den Herren \textsc{Nicolai}, \textsc{Wachter} und \textsc{Gerling}. 30) Beobachtungen und Elemente der Bahn des Kometen von 1811 sowie ein Zusatz zur Theoria motus corporum coelestium von Herrn Professor \textsc{Gauss}. Im \emph{Zusatz} bemerken wir einen Druckfehler. Auf S. 257, Z. 5, muss statt $z$ vielmehr $-z$ gelesen werden. 31) Beobachtungen des Kometen von 1811, die Elemente seiner Bahn und Sternbedeckungen, von Herrn Professor \textsc{Bessel}. 32) Aus einem Schreiben von Herrn Dr. \textsc{Koch} in Danzig. Bemerkungen über einige Fixsterne und Bestimmung der Abweichung der Magnetnadel in Danzig ($13^\circ\,48'$). 33) Beobachtungen des Kometen von 1811 an der königlichen Sternwarte in Berlin. Zu der beigefügten Zeichnung der wahren Kometenbahn ist anzumerken, dass sie keine Projektion ist. --- Unter den kleinen astronomischen Notizen, die den letzten Artikel ausmachen, heben wir nur die Nachricht aus einem Schreiben von Herrn Dr. B\dots\ (\textsc{Brinkley}?) hervor. Er hat mit einem achtfüßigen Kreis die jährliche Parallaxe von $\alpha$ in der Leier zu $2''52$ gefunden, weniger als die Hälfte des oben angeführten Ergebnisses von \textsc{Calandrelli}.

Der Jahrgang 1815 bietet 31 Artikel. Den Anfang machen Bemerkungen von Herrn Professor \textsc{Littrow} zu der Theoria motus corporum coelestium von Herrn Professor \SourceLeaf{556}\textsc{Gauss}. Sie betreffen die Aufgabe, die Bahn eines Planeten aus drei geozentrischen Positionen zu bestimmen. Für den freilich sehr seltenen Fall, dass drei Hypothesen nicht ausreichen, ändert Herr \textsc{Littrow} den Ablauf der Rechnung für die vierte Hypothese geringfügig ab. Mit demselben Kunstgriff, mit dem \textsc{Gauss} die neuen Werte von $P$ und $Q$ bestimmt, berechnet er sogleich $r, r', f, f', f''$. Obwohl dieses Verfahren gelegentlich einen gewissen Vorteil bieten kann, möchten wir es doch nicht uneingeschränkt empfehlen. Im Allgemeinen darf man davon keine genaueren Ergebnisse erwarten als von dem in der Theoria verwendeten Verfahren. Die Abkürzung geht zum Teil wieder verloren, weil man dabei die Lage der Bahnebene nicht mit erhält und es zuweilen etwas heikel sein kann, denselben Kunstgriff auch auf sie sofort anzuwenden. Die darauf folgende Methode desselben Astronomen zur Bestimmung einer Kreisbahn aus zwei geozentrischen Positionen halten wir für sehr zweckmäßig. Wir glauben nur nicht, dass es vorteilhaft ist, sie zur vorläufigen Bestimmung der Abstände bei \emph{Kometenbahnen} anzuwenden. Außerdem finden wir von diesem Verfasser eine einfache Herleitung der Aberration in Länge und Breite und ein paar Reihenausdrücke für die Refraktion, von denen einer nicht hinreichend erklärt ist, um ein Urteil darüber abgeben zu können. 2) Astronomische Beobachtungen aus Pisa, Mailand und Padua (dem XV. Band der Memorie di Matematica e di Fisica della Società Italiana entnommen). 3) Astronomische Nachrichten und Bemerkungen, physische Beobachtung des großen Kometen von 2811, geografische Bestimmungen usw. von Herrn Professor \textsc{Huth} in Dorpat. 4) Beobachtung des großen Kometen von 1811, Berechnung seiner Bahn, Sternbedeckungen usw. von Herrn Professor \textsc{Bessel}. 5) Über die Entdeckung eines neuen Kometen im November 1811, Beobachtungen dieses Kometen und des großen Kometen von 1811 sowie Beobachtung von Pallas, von Herrn Dr. \textsc{Olbers}. 6) Beobachtungen des Kometen von 1807 und die Elemente seiner Bahn, von Herrn \textsc{Niccolo Cacciatore}, Gehilfen an der königlichen Sternwarte in Palermo. Ein sehr wertvoller, bisher auf dem Festland noch. nicht bekannter Nachtrag zur scheinbaren Bahn dieses Kometen. 7) Beobachtete und berechnete Oppositionen von Mars, Vesta, Jupiter und Saturn, Verfinsterungen der Jupitertrabanten und Sternbedeckungen im Jahr 1811, Elemente des Kometen von 1810, Beobachtung des großen Kometen von 1811 und Berechnung seiner Bahn, von Herrn Dr. \textsc{Triesnecker}. 8) Beobachtungen von Uranus, Saturn, Mars und Ceres, von Verfinsterungen der Jupitertrabanten und von Sternbedeckungen sowie des großen Kometen von 1811, an der kaiserlichen Sternwarte in Wilna von Herrn Professor \textsc{Sniadecki} angestellt. 9) Astronomische Nachrichten und Beobachtungen, geografische Ortsbestimmungen usw. von Herrn Professor \textsc{Jabbo Oltmanns}. Beobachtete Sternbedeckungen in Amsterdam und Greenwich, Beobachtungen in Südamerika, Ortsbestimmungen in Spanien und Portugal. 10) Erinnerung an den Kometen von \textsc{Halley}, von Herrn Professor J. W. \textsc{Pfaff} in Nürnberg. Der Verfasser überlässt sich hier der Aussicht auf die vielfältigen Aufschlüsse, die die wiederholt beobachtete Wiederkehr der Kometen einst geben wird. Freilich müssen noch Jahrhunderte vergehen, bevor sich auch nur ein Teil dieser Hoffnungen erfüllen kann. 11) Astronomische Beobachtungen, die 1811 an der königlichen Sternwarte in Berlin angestellt wurden. 12) Beobachtungen über Verfinsterungen der Jupitertrabanten, Mondfinsternisse, Sternbedeckungen, Refraktion, Planetenoppositionen und den großen Kometen von 1811, angestellt an der Prager Sternwarte von Herrn Professor \textsc{David}. 13) Beobachtungen der Pallas und des zweiten Kometen von 1811, Elemente seiner Bahn, Sternbedeckungen usw. von Herrn Professor \textsc{Gauss}. 14) Über das Zusammentreffen der Erde und des Mondes an ein und demselben Ort. Vorausberechnung der Fälle für die Jahre 1812--1815, in denen die Erde einige Stunden später an den Ort gelangt, an dem zuvor der Mond war, und umgekehrt. Lichtenberg hatte bekanntlich einmal die Frage aufgeworfen, ob sich im ersten Fall vielleicht auffällige Wetterveränderungen zeigten. Wir gestehen, dass \SourceLeaf{557}wir (ebenso wie Herr \textsc{Bode} selbst) nichts dergleichen erwarten, und bemerken nur noch, dass die Rechnung selbst auf der unzulässigen Annahme beruht, das Sonnensystem befinde sich in absoluter Ruhe. 15) Beobachtungen der Ceres und der Oppositionen von Uranus und Mars im Jahr 1811, von Herrn \textsc{Derfflinger} an der Sternwarte in Kremsmünster angestellt. 16) Astronomische Bemerkungen von Herrn Dr. \textsc{Lamberti} in Dorpat. 17) Über die scheinbare Bahn des großen Kometen von 1811, mit einer Zeichnung. 18) Über die Bewegung des Doppelsterns 61 im Schwan, von Herrn Professor \textsc{Bessel}. Über diesen bemerkenswerten Gegenstand wurde bereits an mehreren anderen Stellen berichtet, zuerst in unseren Blättern (25. April 1812, 67. Stück [S. 349 d. B.]). 19) Verbesserung der Bestimmung der Polhöhe von Riga, von Herrn Professor \textsc{Sandt}. 20) Astronomische Beobachtungen, die in Paris und Greenwich in den Jahren 1805--1809 angestellt und Herrn Professor \textsc{Oltmanns} mitgeteilt wurden. 21) Nachricht von Herrn Professor \textsc{Fischer} in Berlin über sehr vollkommene Parallelspiegel, die Herr \textsc{Duve}, Mechaniker in Berlin, anfertigt. Da es wichtig ist, dass die Glasspiegel in Reflexionsinstrumenten vollkommen ebene und parallele Flächen haben, sind die Bemühungen von Herrn \textsc{Duve}, die nach Herrn \textsc{Fischer} von Erfolg gekrönt waren, sehr wertvoll. Bemerkenswert ist die hier geschilderte Erfahrung, dass ein solcher Parallelspiegel seine frühere Vollkommenheit verlor, nachdem durch einen Zufall ein kleines Stück davon abgebrochen war. 22) Astronomische Beobachtungen und Bemerkungen von Herrn Joseph \textsc{Bayr} im Kloster Hradisch bei Ollmütz. Geografische Bestimmung von Hradisch und Polhöhe von Troppau. 23) Über den Einfluss der Theorie von \textsc{Dalton} auf die Höhenmessung und die Strahlenbrechung, von Herrn Dr. \textsc{Benzenberg}. Das Urteil über die zwar geistreiche, bisher aber noch zu problematische Hypothese von \textsc{Dalton} gehört vor das Forum des Physikers und nicht vor das des Astronomen. Ob die Erwiderung von Herrn \textsc{Benzenberg} auf die Einwände von Herrn \textsc{Tralles} gegen diese Theorie zulässig ist, müsste die Chemie, so scheint uns, leicht entscheiden können. Vorerst wäre es jedoch verfrüht, jene Hypothese den barometrischen Höhenmessungen als gesicherte Wahrheit zugrunde zu legen. 24) Beiläufige Gedanken über die Oberfläche des Mondes, von Herrn Leutnant \textsc{von Boguslawski}. Obwohl wir auf der Mondoberfläche kein Wasser wahrnehmen, könnte es nach Ansicht von Herrn \textsc{von Boguslawski} dort doch vorhanden sein. Auf der Tagseite würde es von der Sonnenwärme lediglich in unsichtbare gasförmige Dämpfe verwandelt, die nach Sonnenuntergang allmählich, während sie ihre Wärme verlieren, in Nebel, Tau und Eis übergehen. Das augenblickliche Verschwinden der Fixsterne, wenn sie vom dunklen Mondrand bedeckt werden, scheint uns ein wichtiger Einwand gegen diese Hypothese zu sein. 25) Beobachtung von Pallas und Juno, Berechnung ihrer Oppositionen, Elemente der Bahn des Letzteren usw. von Herrn Professor \textsc{Gauss}. Die beiden folgenden Artikel geben die Ephemeriden von Pallas und Juno für 1813 von den Herren \textsc{Nicolai} und \textsc{Wachter}. 28) Noch etwas über den veränderlichen Doppelstern 61 Schwan, Beobachtung des großen Kometen von 1811, Sternbedeckungen und astronomische Nachrichten, von Herrn Professor \textsc{Bessel}. \newpage
29) Astronomische Beobachtungen und Bemerkungen von Herrn Dr. \textsc{Koch} in Danzig. 30) Über den neuen Kometen des Jahres 1812. Der letzte Artikel bringt unter der Überschrift „Vermischte astronomische Nachrichten“ weitere Beobachtungen des großen Kometen von 1811 in Petersburg durch Herrn \textsc{Schubert}, die Preise der astronomischen Instrumente, die Herr \textsc{Baumann} in Stuttgart anfertigt, und verschiedene andere astronomische Notizen. --- Wir schließen diese Anzeige des an Reichhaltigkeit und Interesse stets gleichbleibenden Jahrbuchs mit dem Wunsch, dass der würdige Herausgeber auch künftig bei seinem mühsamen Unternehmen tatkräftig unterstützt werde und uns noch lange jedes Jahr mit der Fortsetzung dieses Jahrbuchs erfreuen möge.
